\documentclass[11pt]{article}

\usepackage[utf8]{inputenc}
\usepackage[T1]{fontenc}
\usepackage{lmodern}
\usepackage{geometry}
\usepackage{amsmath,amssymb,amsfonts,amsthm,mathtools}
\usepackage{bm}
\usepackage{enumitem}
\usepackage{microtype}
\usepackage{hyperref}
\usepackage{titlesec}
\usepackage{booktabs}
\usepackage{array}
\usepackage{setspace}
\usepackage{mathrsfs}
\usepackage{physics}

\geometry{margin=1in}
\hypersetup{colorlinks=true, linkcolor=black, urlcolor=blue, citecolor=black}
\setlist[enumerate]{topsep=0.4em,itemsep=0.35em}
\setlist[itemize]{topsep=0.3em,itemsep=0.25em}
\titleformat{\section}{\large\bfseries}{\thesection.}{0.5em}{}
\titleformat{\subsection}{\normalsize\bfseries}{\thesubsection.}{0.5em}{}
\setstretch{1.06}

\newcommand{\Aether}{\AE{}ther}
\newcommand{\Aflow}{\AE{}ther-flow}
\newcommand{\TheoryName}{\Aflow{} Interpretation of Relativity}
\newcommand{\TheoryProgram}{\Aether{} / \Aflow{} framework}

\newcommand{\CompletedMetric}{g^{\sharp}}
\newcommand{\MatterSpace}{\Psi}

\newcommand{\Car}{\mathrm{car}}
\newcommand{\Cur}{\mathrm{cur}}
\newcommand{\Hom}{\mathrm{hom}}

\newcommand{\ForSpace}[1]{\mathcal I_{#1}^{\mathrm{for}}}
\newcommand{\PrimShell}{\mathcal S_{\mathrm{prim}}}
\newcommand{\Reduction}[1]{\mathcal R_{#1}}
\newcommand{\CurrentCarrier}[1]{\mathfrak C_{#1}^{\mathrm{cur}}}
\newcommand{\EqCarrier}[1]{\mathfrak C_{#1}^{\mathrm{eq}}}
\newcommand{\MetricOut}{\mathscr G_{\mathrm{out}}}
\newcommand{\MatterOut}{\mathscr T_{\mathrm{out}}}

\newcommand{\DataSpace}[1]{\mathcal Y_{#1}^{\mathrm{obs}}}
\newcommand{\AmbientTarget}[1]{E_{#1}^{\mathrm{amb}}}
\newcommand{\EqnTarget}[1]{Q_{#1}^{\mathrm{eqn}}}

\newcommand{\FoundSpace}[1]{\mathcal F_{#1}^{\mathrm{fdn}}}
\newcommand{\GroundSpace}[1]{\mathcal G_{#1}^{\mathrm{grd}}}
\newcommand{\OrigSpace}[1]{\mathcal O_{#1}^{\mathrm{orig}}}
\newcommand{\PrecSpace}[1]{\mathcal P_{#1}^{\mathrm{prec}}}
\newcommand{\LiftSpace}[1]{\mathcal L_{#1}^{\mathrm{lift}}}
\newcommand{\SubSpace}[1]{\mathcal M_{#1}^{\mathrm{sub}}}
\newcommand{\SubTarget}[1]{E_{#1}^{\mathrm{sub}}}
\newcommand{\SubPair}[1]{\mathfrak s_{#1}^{\mathrm{sub}}}

\newcommand{\HiddenSpace}[1]{\mathcal H_{#1}^{\mathrm{hid}}}
\newcommand{\HiddenBody}[1]{D_{#1}^{\mathrm{hid}}}

\newcommand{\SeedSpace}[1]{\mathcal S_{#1}^{\mathrm{seed}}}
\newcommand{\SeedBody}[1]{\mathcal B_{#1}^{\mathrm{seed}}}
\newcommand{\SeedProj}[1]{\Pi_{#1}^{\mathrm{seed}}}
\newcommand{\SeedFoundMin}[1]{\mathfrak f_{#1}^{\mathrm{seed,fdn}}}
\newcommand{\SeedGroundMin}[1]{\mathfrak f_{#1}^{\mathrm{seed,grd}}}
\newcommand{\SeedOrigMin}[1]{\mathfrak f_{#1}^{\mathrm{seed,orig}}}
\newcommand{\SeedPrecMin}[1]{\mathfrak f_{#1}^{\mathrm{seed,prec}}}
\newcommand{\SeedLiftMin}[1]{\mathfrak f_{#1}^{\mathrm{seed,lift}}}
\newcommand{\SeedSubMin}[1]{\mathfrak m_{#1}^{\mathrm{seed}}}
\newcommand{\SeedSection}[1]{\Gamma_{#1}^{\mathrm{seed}}}
\newcommand{\SeedFunctional}[1]{\mathscr W_{#1}^{\mathrm{seed}}}
\newcommand{\SeedShellSet}[1]{\mathcal Z_{#1}^{\mathrm{seed}}}
\newcommand{\SeedZeroCore}[1]{\mathcal Z_{#1}^{0,\mathrm{seed}}}
\newcommand{\SeedSplit}[1]{\Phi_{#1}^{\mathrm{seed}}}
\newcommand{\SeedCharge}[1]{\mathcal Q_{#1}^{\mathrm{seed}}}
\newcommand{\SeedEmbed}[1]{\zeta_{#1}^{\mathrm{seed}}}
\newcommand{\SeedKernelCh}[1]{\widetilde K_{#1}^{0,\mathrm{seed}}}
\newcommand{\SeedStateComp}[1]{P_{#1}^{\mathrm{seed,co}}}

\newcommand{\RootNucleus}[1]{\mathfrak N_{#1}^{\mathrm{rt}}}
\newcommand{\RootProj}[1]{\Pi_{#1}^{\mathrm{rt}}}
\newcommand{\RootFoundMin}[1]{\mathfrak f_{#1}^{\mathrm{rt,fdn}}}
\newcommand{\RootGroundMin}[1]{\mathfrak f_{#1}^{\mathrm{rt,grd}}}
\newcommand{\RootOrigMin}[1]{\mathfrak f_{#1}^{\mathrm{rt,orig}}}
\newcommand{\RootPrecMin}[1]{\mathfrak f_{#1}^{\mathrm{rt,prec}}}
\newcommand{\RootLiftMin}[1]{\mathfrak f_{#1}^{\mathrm{rt,lift}}}
\newcommand{\RootSubMin}[1]{\mathfrak m_{#1}^{\mathrm{rt}}}
\newcommand{\RootShellSet}[1]{\mathcal Z_{#1}^{\mathrm{rt}}}
\newcommand{\RootZeroCore}[1]{\mathcal Z_{#1}^{0,\mathrm{rt}}}
\newcommand{\RootSplit}[1]{\Phi_{#1}^{\mathrm{rt}}}
\newcommand{\RootCharge}[1]{\mathcal Q_{#1}^{\mathrm{rt}}}
\newcommand{\RootEmbed}[1]{\zeta_{#1}^{\mathrm{rt}}}
\newcommand{\RootKernelCh}[1]{\widetilde K_{#1}^{0,\mathrm{rt}}}
\newcommand{\RootStateComp}[1]{P_{#1}^{\mathrm{rt,co}}}

\newcommand{\SeedMapFound}[1]{\mathcal E_{#1}^{\mathrm{seed,fdn}}}
\newcommand{\SeedMapGround}[1]{\mathcal E_{#1}^{\mathrm{seed,grd}}}
\newcommand{\SeedMapOrig}[1]{\mathcal E_{#1}^{\mathrm{seed,orig}}}
\newcommand{\SeedMapPrec}[1]{\mathcal E_{#1}^{\mathrm{seed,prec}}}
\newcommand{\SeedMapLift}[1]{\mathcal E_{#1}^{\mathrm{seed,lift}}}
\newcommand{\SeedMapSub}[1]{\mathcal E_{#1}^{\mathrm{seed,sub}}}
\newcommand{\SeedMapShell}[1]{\mathcal E_{#1}^{\mathrm{seed,sh}}}
\newcommand{\SeedMapZero}[1]{\mathcal E_{#1}^{\mathrm{seed},0}}

\newtheorem{definition}{Definition}
\newtheorem{proposition}{Proposition}
\newtheorem{remark}{Remark}
\newtheorem{corollary}{Corollary}
\newtheorem{theorem}{Theorem}

\title{\textbf{The \TheoryName{}}\\[0.4em]
\large Primitive-Reservoir Observer-Reduction\\
\large Minimal Root Exact-Shell Transport Nucleus\\
\large from Nucleus-Generating Substrate Seed Dynamics}
\author{}
\date{}

\begin{document}

\maketitle

\begin{abstract}
The current primitive-reservoir observer-reduction frontier is the minimal root exact-shell transport nucleus. The previous collapse theorem already spent the remaining internal compression move available at root scope: the larger benchmark-facing root core carries no extra benchmark-relevant freedom beyond that smaller nucleus. After that result, the next honest benchmark-facing task is no longer another same-output root, foundation, or ground relay. It is either to derive or justify the nucleus from more primitive substrate structure, or to prove a genuinely smaller theorem-level collapse strictly inside the nucleus.

The present note takes the first route. For each sector it introduces \emph{nucleus-generating substrate seed dynamics}: a compact convex seed body over the recorded foundation fiber, unique seed minimizers on the fixed foundation, ground, origin, precursor, lift, and substrate fibers, an exact seed restriction section, an exact seed shell, a closed embedded seed zero core whose hidden split saturates that shell, a hidden charge, an observer-compatible exact zero-response realization, and the fixed-data solvability / coercive package. The minimal root exact-shell transport nucleus is then recovered canonically as the exact-shell transport image of that deeper seed dynamics.

The benchmark boundary remains fixed throughout. The one operative metric is still exactly \(\CompletedMetric\), the ordinary matter route is unchanged, there is no second operative metric, and the low-energy matter law is not enlarged. The positive result is therefore narrow but benchmark facing: the minimal root exact-shell transport nucleus is no longer left as an unexplained primitive stopping point on the active line. The remaining honest burden moves one layer deeper again, to the seed dynamics themselves or to a sharper collapse inside their observer-relevant core.
\end{abstract}

\section{Introduction}

The active derivational program of \TheoryName{} has now reached a cleaner burden than another same-output relay. The current root-nucleus collapse theorem already proved that the benchmark-facing root data collapse to a smaller minimal root exact-shell transport nucleus. The accompanying routing note then fixed the consequence of that theorem: the next honest benchmark-facing manuscript must either derive or justify that nucleus from more primitive substrate structure, or prove a genuinely smaller collapse strictly inside it. Reopening the former larger root core, re-running another same-output root / foundation / ground relay, or returning to stationary-reduction, stationary-Hessian, spectral, or selector layers would no longer address the live burden.

The present note takes the first route. Its positive claim is not that final substrate microphysics has been identified, and it is not that the benchmark exact-relativistic package has been changed. Its claim is narrower and benchmark facing. The current root-nucleus datum is shown to arise canonically from a more primitive seed-level substrate dynamics whose exact-shell transport already contains the observer-relevant structure of the nucleus. In that sense the nucleus is no longer the primitive place where the active line has to stop.

Two features matter here. First, the theorem targets the current minimal nucleus itself rather than the former larger root core, so it does not spend the next move on an already-compressed presentation. Second, the deeper seed package is not merely another renaming of the nucleus. It adds an explicit seed state space, compact seed body, restriction section, and exact-shell dynamics below the nucleus. The shell and zero-core data that the current frontier treats as benchmark-facing are therefore transported exact-shell outputs of a deeper substrate package rather than unexplained primitive data.

\paragraph{Framework and claim boundary.}
\TheoryProgram{} denotes the broader \Aether{} / \Aflow{} framework built on the claims that the \Aether{} is the underlying four-dimensional substrate of reality, the \Aflow{} is its intrinsic ordered motion, observed three-dimensional space is the local experiential slice of that deeper substrate, S-time is the experienced order of change arising from matter, light, and the \Aflow{}, and observed expansion is the three-dimensional appearance of deeper four-dimensional ordered motion. Within that framework, \TheoryName{} denotes the exact-closure theory in which the effective gravitational dynamics are adopted to be exactly Einsteinian with universal matter coupling. In the active sequence, \emph{adoption} means use of that established relativistic dynamics without claiming substrate derivation, while \emph{derivation} is reserved for a first-principles recovery from explicit substrate variables.

The present note stays strictly inside that boundary. It does not alter the benchmark exact-closure package. It does not introduce a second operative metric or a new low-energy matter law. It only proves that the current minimal root exact-shell transport nucleus can itself be generated from a more primitive seed-level substrate dynamics.

\section{Fixed Inputs}

\begin{definition}[Recorded primitive continuation data]
\label{def:fixed}
Fix the following continuation data from the active primitive-reservoir observer-reduction line.
\begin{enumerate}
    \item For each sector label \(\sigma\in\{\Car,\Cur,\Hom\}\), a primitive forcing shell
    \[
    \ForSpace{\sigma},
    \]
    a visible reduction map
    \[
    \Reduction{\sigma}:\ForSpace{\sigma}\to X_{\sigma},
    \]
    and shell-level current and equilibrium carriers
    \[
    \CurrentCarrier{\sigma}:\ForSpace{\sigma}\to Z_{\sigma}^{\mathrm{cur}},
    \qquad
    \EqCarrier{\sigma}:\ForSpace{\sigma}\to Z_{\sigma}^{\mathrm{eq}}.
    \]
    \item The product shell
    \[
    \PrimShell
    \equiv
    \ForSpace{\Car}\times \ForSpace{\Cur}\times \ForSpace{\Hom}.
    \]
    \item The recorded metric and ordinary-matter outputs
    \[
    \MetricOut:\PrimShell\to\{\text{observer metrics}\},
    \qquad
    \MatterOut:\PrimShell\times\MatterSpace\to\{\text{observer stress tensors}\},
    \]
    satisfying
    \begin{equation}
    \MetricOut(\mathbf w)=\CompletedMetric,
    \qquad
    \MatterOut(\mathbf w,\psi)
    =
    T_{\mu\nu}^{\mathrm{matter}}[\CompletedMetric,\psi]
    \label{eq:fixedoutputs}
    \end{equation}
    for every \(\mathbf w\in\PrimShell\) and every matter configuration \(\psi\in\MatterSpace\).
\end{enumerate}
\end{definition}

\begin{definition}[Recorded lower-fiber ladder]
\label{def:ladder}
Fix \(\sigma\in\{\Car,\Cur,\Hom\}\). The active line already fixes the lower-fiber ladder
\[
\FoundSpace{\sigma},
\qquad
\GroundSpace{\sigma},
\qquad
\OrigSpace{\sigma},
\qquad
\PrecSpace{\sigma},
\qquad
\LiftSpace{\sigma},
\qquad
\SubSpace{\sigma},
\]
together with the common substrate target \(\SubTarget{\sigma}\), the positive-definite pairing
\[
\SubPair{\sigma}:
\SubTarget{\sigma}\times\SubTarget{\sigma}\to\mathbb R,
\]
and the hidden equation target \(\EqnTarget{\sigma}\). The current root-nucleus frontier uses exactly this recorded ladder when it names the projection and the six minimizer images that remain benchmark facing at root scope.
\end{definition}

\begin{remark}
Definitions \ref{def:fixed} and \ref{def:ladder} are not reopened here. The primitive continuation data, the one operative metric, the ordinary matter route, and the recorded lower-fiber ladder remain fixed inputs. The present note addresses only whether the current minimal root exact-shell transport nucleus can itself be generated from more primitive substrate dynamics.
\end{remark}

\begin{remark}[Downstream fixed fact]
The already-recorded current frontier theorem proves that once a minimal root exact-shell transport nucleus is fixed, the larger benchmark-facing root core is canonically induced and carries no extra benchmark-relevant freedom. The present note uses that downstream fact only as a routing consequence: after the induced nucleus is obtained, no second larger root package needs to be postulated anew.
\end{remark}

\section{Nucleus-Generating Substrate Seed Dynamics}

\begin{definition}[Nucleus-generating substrate seed dynamics]
\label{def:seed}
Fix \(\sigma\in\{\Car,\Cur,\Hom\}\). A \emph{nucleus-generating substrate seed dynamics package} for sector \(\sigma\) consists of the following data.
\begin{enumerate}
    \item A finite-dimensional Euclidean seed state space \(\SeedSpace{\sigma}\), a compact convex seed body
    \[
    \SeedBody{\sigma}\subset\SeedSpace{\sigma}
    \]
    with nonempty interior, an affine surjection
    \[
    \SeedProj{\sigma}:\SeedSpace{\sigma}\to\FoundSpace{\sigma},
    \]
    and smooth seed fiber minimizers
    \[
    \SeedFoundMin{\sigma},\quad
    \SeedGroundMin{\sigma},\quad
    \SeedOrigMin{\sigma},\quad
    \SeedPrecMin{\sigma},\quad
    \SeedLiftMin{\sigma},\quad
    \SeedSubMin{\sigma}
    \]
    on the fixed foundation, ground, origin, precursor, lift, and substrate fibers of the recorded ladder.
    \item A smooth seed restriction section
    \[
    \SeedSection{\sigma}:\SeedSpace{\sigma}\to\SubTarget{\sigma},
    \]
    the seed functional
    \[
    \SeedFunctional{\sigma}(s)
    \equiv
    \SubPair{\sigma}\bigl(\SeedSection{\sigma}(s),\SeedSection{\sigma}(s)\bigr),
    \]
    and the exact seed shell
    \[
    \SeedShellSet{\sigma}
    \equiv
    \bigl(\SeedSection{\sigma}\bigr)^{-1}(0)\cap\SeedBody{\sigma}.
    \]
    \item A closed embedded seed zero core
    \[
    \SeedZeroCore{\sigma}\subset\SeedShellSet{\sigma},
    \]
    a finite-dimensional hidden fiber space \(\HiddenSpace{\sigma}\) with compact hidden body \(\HiddenBody{\sigma}\subset\HiddenSpace{\sigma}\), and a split diffeomorphism
    \[
    \SeedSplit{\sigma}:
    \SeedZeroCore{\sigma}\times\HiddenBody{\sigma}
    \xrightarrow{\cong}
    \SeedShellSet{\sigma}.
    \]
    \item A hidden charge
    \[
    \SeedCharge{\sigma}:\HiddenSpace{\sigma}\to\EqnTarget{\sigma},
    \]
    a smooth observer-compatible exact seed zero-response realization
    \[
    \SeedEmbed{\sigma}:\ForSpace{\sigma}\hookrightarrow\SeedZeroCore{\sigma},
    \]
    and the fixed-data solvability / coercive package on \(\SeedZeroCore{\sigma}\), recorded through the charted kernel \(\SeedKernelCh{\sigma}\) and orthogonal projector \(\SeedStateComp{\sigma}\).
\end{enumerate}
The package is \emph{benchmark faithful} if it preserves the benchmark outputs \eqref{eq:fixedoutputs}, introduces no second operative metric, introduces no enlarged low-energy matter law, and is presented as a deeper substrate-side source for the current minimal root exact-shell transport nucleus rather than as a replacement benchmark theory.
\end{definition}

\begin{remark}
Definition \ref{def:seed} is intentionally narrower than another full root relay. It does not rebuild the larger former root core as an independent datum. It only introduces the deeper seed state space and exact-shell dynamics needed to generate the current minimal nucleus.
\end{remark}

\section{The Induced Root Nucleus}

\begin{definition}[Seed-induced minimal root exact-shell transport nucleus]
\label{def:induced}
Fix a benchmark-faithful seed package of Definition \ref{def:seed}. Its \emph{seed-induced minimal root exact-shell transport nucleus} \(\RootNucleus{\sigma}\) is the tuple consisting of:
\begin{enumerate}
    \item the induced root projection and the six induced minimizer images
    \[
    \RootProj{\sigma}\equiv\SeedProj{\sigma},
    \qquad
    \RootFoundMin{\sigma}\equiv\SeedFoundMin{\sigma},
    \qquad
    \RootGroundMin{\sigma}\equiv\SeedGroundMin{\sigma},
    \]
    \[
    \RootOrigMin{\sigma}\equiv\SeedOrigMin{\sigma},
    \qquad
    \RootPrecMin{\sigma}\equiv\SeedPrecMin{\sigma},
    \qquad
    \RootLiftMin{\sigma}\equiv\SeedLiftMin{\sigma},
    \qquad
    \RootSubMin{\sigma}\equiv\SeedSubMin{\sigma};
    \]
    \item the induced exact root shell and embedded root zero core
    \[
    \RootShellSet{\sigma}\equiv\SeedShellSet{\sigma},
    \qquad
    \RootZeroCore{\sigma}\equiv\SeedZeroCore{\sigma};
    \]
    \item the induced split diffeomorphism
    \[
    \RootSplit{\sigma}\equiv\SeedSplit{\sigma}:
    \RootZeroCore{\sigma}\times\HiddenBody{\sigma}
    \xrightarrow{\cong}
    \RootShellSet{\sigma};
    \]
    \item the induced hidden charge, exact zero-response realization, and solvability / coercive package
    \[
    \RootCharge{\sigma}\equiv\SeedCharge{\sigma},
    \qquad
    \RootEmbed{\sigma}\equiv\SeedEmbed{\sigma},
    \]
    together with
    \[
    \RootKernelCh{\sigma}\equiv\SeedKernelCh{\sigma},
    \qquad
    \RootStateComp{\sigma}\equiv\SeedStateComp{\sigma}.
    \]
\end{enumerate}
\end{definition}

\begin{definition}[Canonical seed comparison maps]
\label{def:equiv}
Let \({}^{(1)}\!\RootNucleus{\sigma}\) and \({}^{(2)}\!\RootNucleus{\sigma}\) be the nuclei induced by two benchmark-faithful seed packages over the same fixed lower-fiber ladder. Their canonical comparison maps are defined by the common lower data they induce:
\begin{align}
\SeedMapFound{\sigma}\bigl({}^{(1)}\!\SeedFoundMin{\sigma}(f)\bigr)
&\equiv
{}^{(2)}\!\SeedFoundMin{\sigma}(f),
\\
\SeedMapGround{\sigma}\bigl({}^{(1)}\!\SeedGroundMin{\sigma}(g)\bigr)
&\equiv
{}^{(2)}\!\SeedGroundMin{\sigma}(g),
\\
\SeedMapOrig{\sigma}\bigl({}^{(1)}\!\SeedOrigMin{\sigma}(o)\bigr)
&\equiv
{}^{(2)}\!\SeedOrigMin{\sigma}(o),
\\
\SeedMapPrec{\sigma}\bigl({}^{(1)}\!\SeedPrecMin{\sigma}(p)\bigr)
&\equiv
{}^{(2)}\!\SeedPrecMin{\sigma}(p),
\\
\SeedMapLift{\sigma}\bigl({}^{(1)}\!\SeedLiftMin{\sigma}(\ell)\bigr)
&\equiv
{}^{(2)}\!\SeedLiftMin{\sigma}(\ell),
\\
\SeedMapSub{\sigma}\bigl({}^{(1)}\!\SeedSubMin{\sigma}(m)\bigr)
&\equiv
{}^{(2)}\!\SeedSubMin{\sigma}(m),
\\
\SeedMapShell{\sigma}
&\equiv
\bigl({}^{(2)}\!\SeedProj{\sigma}\big|_{{}^{(2)}\!\SeedShellSet{\sigma}}\bigr)^{-1}
\circ
{}^{(1)}\!\SeedProj{\sigma}\big|_{{}^{(1)}\!\SeedShellSet{\sigma}},
\\
\SeedMapZero{\sigma}
&\equiv
\bigl({}^{(2)}\!\SeedProj{\sigma}\big|_{{}^{(2)}\!\SeedZeroCore{\sigma}}\bigr)^{-1}
\circ
{}^{(1)}\!\SeedProj{\sigma}\big|_{{}^{(1)}\!\SeedZeroCore{\sigma}}.
\end{align}
The two induced nuclei are \emph{seed-nuclearly equivalent} if these comparison maps transport every displayed structure in Definition \ref{def:induced} identically and intertwine the split maps via
\[
\SeedMapShell{\sigma}\circ{}^{(1)}\!\SeedSplit{\sigma}
=
{}^{(2)}\!\SeedSplit{\sigma}\circ
\bigl(\SeedMapZero{\sigma}\times\mathrm{id}_{\HiddenBody{\sigma}}\bigr).
\]
\end{definition}

\begin{proposition}[The shell is generated, not postulated]
\label{prop:shell}
Fix a benchmark-faithful seed package. Then
\[
\RootShellSet{\sigma}
=
\operatorname{im}\RootSplit{\sigma}
\]
for the induced nucleus of Definition \ref{def:induced}.
\end{proposition}

\begin{proof}
By Definition \ref{def:seed}, the split map \(\SeedSplit{\sigma}\) is a diffeomorphism from \(\SeedZeroCore{\sigma}\times\HiddenBody{\sigma}\) onto the full seed shell \(\SeedShellSet{\sigma}\). Definition \ref{def:induced} identifies \(\RootSplit{\sigma}\) with \(\SeedSplit{\sigma}\) and \(\RootShellSet{\sigma}\) with \(\SeedShellSet{\sigma}\). Therefore the induced shell is exactly the image of the induced split.
\end{proof}

\begin{proposition}[Seed dynamics canonically induce the current root-nucleus data]
\label{prop:induce}
Fix a benchmark-faithful seed package. Then the data of Definition \ref{def:induced} satisfy exactly the shape of the current minimal root exact-shell transport nucleus.
\end{proposition}

\begin{proof}
The induced projection and the six induced minimizer images are supplied directly by Definition \ref{def:seed}(1). The induced exact shell and embedded zero core come from Definition \ref{def:seed}(2)--(3). Proposition \ref{prop:shell} shows that, inside this deeper class, the shell is generated by the split rather than added as an extra primitive datum. The hidden charge, the exact zero-response realization, and the solvability / coercive package are provided by Definition \ref{def:seed}(4). These are precisely the constituents retained by the current root-nucleus frontier.
\end{proof}

\begin{proposition}[The benchmark package remains fixed]
\label{prop:boundary}
For every benchmark-faithful seed package, the operative metric remains exactly \(\CompletedMetric\), the ordinary matter route remains unchanged, no second operative metric is introduced, and the low-energy matter law is not enlarged.
\end{proposition}

\begin{proof}
This is part of the benchmark-faithfulness requirement in Definition \ref{def:seed}. Equation \eqref{eq:fixedoutputs} remains unchanged on the full primitive forcing shell, so the induced nucleus does not alter the benchmark exact-closure package.
\end{proof}

\section{Main Theorem}

\begin{theorem}[Minimal root exact-shell transport nucleus from nucleus-generating substrate seed dynamics]
\label{thm:main}
Fix the primitive continuation data of Definition \ref{def:fixed} and the recorded lower-fiber ladder of Definition \ref{def:ladder}. Then, for each sector \(\sigma\in\{\Car,\Cur,\Hom\}\), every benchmark-faithful nucleus-generating substrate seed dynamics package canonically determines a minimal root exact-shell transport nucleus with the same benchmark one-metric outputs. More precisely:
\begin{enumerate}
    \item every benchmark-faithful seed package canonically determines an induced nucleus \(\RootNucleus{\sigma}\) in the sense of Definition \ref{def:induced};
    \item if two benchmark-faithful seed packages are seed-nuclearly equivalent, then their induced nuclei are benchmark-equivalent at current root-nucleus scope;
    \item the current root-nucleus frontier is therefore no longer an unexplained primitive stopping point on the active line, but the exact-shell transport image of a more primitive seed dynamics;
    \item by the already-recorded current frontier theorem, the full benchmark-facing exact-shell root core is then recovered downstream from that induced nucleus without any new larger root datum.
\end{enumerate}
\end{theorem}

\begin{proof}
Item (1) is Proposition \ref{prop:induce}. For item (2), the canonical comparison maps of Definition \ref{def:equiv} transport the induced projection, all six induced minimizer images, the induced shell and zero core, the induced split, the hidden charge, the exact zero-response realization, and the solvability / coercive package identically. Hence the two induced nuclei agree at exactly the benchmark-facing level retained by the current frontier.

Item (3) then follows immediately. Once the seed dynamics are fixed, the current nucleus is no longer introduced as an independent primitive tuple. It is the transported exact-shell output of the deeper seed state space, seed body, and seed restriction dynamics. This is a genuine change in benchmark-facing status because it removes the nucleus itself as the unexplained stopping point rather than merely restating the former larger root package.

Item (4) is the downstream routing consequence already fixed on the active line: the present theorem produces the current minimal nucleus, and the already-recorded frontier theorem then recovers the full benchmark-facing root core from that nucleus. No second larger root package is postulated here, and no new benchmark freedom is added there.
\end{proof}

\begin{corollary}[What must be done next]
\label{cor:next}
After Theorem \ref{thm:main}, the next honest benchmark-facing manuscript must either:
\begin{enumerate}
    \item derive or justify the nucleus-generating substrate seed dynamics from still more primitive substrate structure in a way that changes benchmark-facing status;
    \item identify a genuinely smaller theorem-level observer-relevant core strictly inside the seed dynamics and prove a sharper collapse to it;
    \item do either of the previous two while preserving the same benchmark one-metric package and the same claim boundary.
\end{enumerate}
Another same-output restatement of the minimal root exact-shell transport nucleus, another same-output root / foundation / ground relay, a reopening of stationary-reduction, stationary-Hessian, spectral, or selector layers, or any move that introduces a second operative metric, enlarges the ordinary matter law, or uses stronger promotion language does not satisfy the present burden.
\end{corollary}

\begin{proof}
Theorem \ref{thm:main} removes the current minimal root exact-shell transport nucleus as the unexplained stopping point on the active line. Therefore another manuscript that merely restates that same nucleus, or merely returns to an already-screened larger relay, does not advance the benchmark-facing status of the project. What remains open is either to go one honest layer deeper again, to the seed dynamics themselves, or to discover a smaller theorem-level observer-relevant core inside the seed package.
\end{proof}

\section{Conclusion}

The positive result of this note is narrow but benchmark facing. The current root-nucleus collapse theorem had already shown that the larger benchmark-facing root core contains no extra benchmark-relevant freedom beyond a minimal root exact-shell transport nucleus. The present theorem then removes the next unexplained stopping point by showing that this nucleus itself is the exact-shell transport image of a more primitive nucleus-generating substrate seed dynamics.

The scope remains disciplined. The benchmark package is unchanged: the one operative metric is still exactly \(\CompletedMetric\), the ordinary matter route is unchanged, there is no second operative metric, and the low-energy matter law is not enlarged. Nothing here upgrades adoption to completed derivation or licenses stronger promotion language.

The open burden therefore moves one honest step further again. The next benchmark-facing manuscript should derive or justify the seed dynamics from still more primitive substrate structure, unless the sharper honest gain available instead is a still smaller theorem-level collapse inside the observer-relevant seed package. Until then, the current result should be read for what it is: a deeper nucleus-scope derivation on the active line of \TheoryName{}, not a claim that the full first-principles program is finished.

\end{document}
